\section{K-024: the unique maximal dyadic hook-chain component}

Use the auxiliary representation and weighted Hilbert series $B_n(t)$
of K-023.  Realize $\mathbb S_{(2,1)}W$ as the degree-three free
Lie submodule of $W^{\otimes3}$, with $[a,b]=ab-ba$.
Put $M=2^{n-1}$ and
\[
 z_n=(e_n,F_n),\qquad
            F_n=\sum_{i=1}^{n-1}[e_i,[e_i,e_{i+1}]].
\]

\begin{theorem}[K-024]
For every $n\ge2$, the following assertions hold.
\begin{enumerate}
\item The $SL_n$ orbit of $z_n$ is closed, its full $GL_n$
stabilizer is cyclic of order $M$, and determinant is faithful
on that stabilizer.  The $SL_n$ stabilizer is trivial.
\item In the affine quotient of $\mathcal U_n$ by $SL_n$, the
nonzero fixed locus of the grading subgroup $\mu_M$ is a single
grading orbit.  Its closure is a density-one curve, smooth at
every nonzero point of the ambient quotient.
\item At each primitive $M$th root $\xi$, the series $B_n(t)$
has a simple pole with
\[
                  C_n(\xi)=\lim_{t\to\xi}(1-t^M)B_n(t)
                       \in\mathbb Z[\xi,1/2]^\times.
                                                               \tag{K-024a}
\]
\item If an integer $d$ has two-part at least $M$ and $d\ne M$,
then $B_n(t)$ is regular at every primitive $d$th root.
In particular there are no poles at $M3^r$ for $r>0$.
\end{enumerate}
\end{theorem}

\begin{proof}
For $1\le i<n$ put
\[
             q_i=\frac{1-(-1/2)^i}{3},\qquad
             q_n=1-q_{n-1},\qquad
 a_i=2e_i+e_{i+1}\ (i<n),\quad a_n=e_n.
\]
All $q_i$ are positive, and
$\sum_i q_i a_i=\mathbf1$: explicitly $2q_1=1$,
$q_{i-1}+2q_i=1$, and $q_{n-1}+q_n=1$.
The Lie words $[e_i,[e_i,e_{i+1}]]$ have squared norm six in
the tensor-product Hermitian metric.  Take the coefficient of
this word to be $\sqrt{q_i/6}$ and that of the vector $e_n$ to
be $\sqrt{q_n}$.  Their diagonal moment sum is $\mathbf1$.
Off-diagonal moments vanish: distinct displayed cubic weights
do not differ by a root $e_a-e_b$, and the linear and cubic
summands are orthogonal.  The resulting representative is
critical for $SL_n$, so its orbit is closed by the Kempf--Ness
criterion used in P-F60 and K-018.

The $n$ displayed weight columns are independent: their matrix
is triangular with determinant $2^{n-1}$.  Consequently the
diagonal torus acts transitively on their nonzero coefficient
choices.  Our critical representative is diagonally equivalent
to $z_n$, and the $SL_n$ orbit of $z_n$ is also closed.

Choose a primitive $M$th root $\zeta$ and write
\[
 w_i=(-2)^{i-1}\pmod M,\qquad
       g=\operatorname{diag}(\zeta^{w_1},\ldots,\zeta^{w_n}).
                                                               \tag{K-024b}
\]
These weights are distinct, $w_n=0$, and
$2w_i+w_{i+1}=0$.  Thus $g$ fixes $z_n$.
Its determinant exponent $\sum_iw_i$ is odd.  The exact
finite-group computation below identifies the full stabilizer.

Let $L=\operatorname{Stab}_{GL_n}(z_n)$.  The independent all-$n$
Lie calculation R-NEW-004 proves that $L$ is finite.  Since $g$
has simple spectrum, its centralizer in $L$ is diagonal.
A diagonal element of $L$ satisfies
\[
                     d_n=1,\qquad d_i^2d_{i+1}=1\quad(i<n).
\]
These equations give precisely the cyclic group $\langle g\rangle$
of order $M$.  Character column orthogonality therefore gives
\[
                  \sum_{\rho\in\operatorname{Irr}(L)}
                              |\rho(g)|^2=M.
\]
The first $M$ powers of the determinant character are distinct
linear characters, so they exhaust this sum.  There are exactly
$M$ linear characters, they are all powers of determinant,
and every nonlinear irreducible character vanishes at $g$.
In particular $L^{\rm ab}\simeq\mu_M$ via determinant.

Set $K=\ker(\det:L\to\mathbb C^*)$.  A zero trace of a
representation of $g$ is equivalent to equality of the
multiplicities at each antipodal pair of eigenvalues: this follows
from the basis $1,\zeta,\ldots,\zeta^{M/2-1}$ of
$\mathbb Q(\zeta)$ and $\zeta^{M/2}=-1$.
The spectrum in (K-024b) has exactly one antipodal pair, on
\[
                           P=\langle e_{n-1},e_n\rangle.
\]
Hence every nonlinear irreducible constituent of $W$ is contained
in this two-plane.  The complementary eigenspaces are linear
characters of $L$, so $K$ fixes them pointwise.  The plane $P$
is $L$-stable (either one nonlinear constituent or two linear
ones), and $K$ acts on it with determinant one.  It fixes the
nonzero vector $e_n$.  A finite-order matrix in $SL_2(\mathbb C)$
fixing a nonzero vector has both eigenvalues one and is the
identity.  Thus $K$ is trivial on $P$ and on its complement.
The representation $L\subset GL(W)$ is faithful, so $K=1$.
Therefore $L=\langle g\rangle$, and the $SL_n$ stabilizer is
trivial.  This argument also covers $n=2$.

We next classify the entire nonzero fixed locus.  Let $y$ be a
closed $SL_n$ orbit whose quotient image is nonzero and fixed by
a primitive grading element of order $M$.  There is an exact
fixing element in $GL_n$ whose determinant has order $M$.
It may be replaced by one of finite two-power order with the
same primitive determinant order.  For completeness, take its
semisimple part and the Zariski closure of the cyclic group it
generates.  A nonzero positive-degree invariant at $y$ makes
determinant finite on this diagonalizable group.  Determinant
is therefore trivial on its connected torus.  A splitting of
its finitely generated character module lifts the desired
finite two-primary component to a finite two-power element.
The odd-primary part can be discarded.  This is the same finite
lifting argument used in K-009; it imposes no bound on the lift
order before classification.

Diagonalize this lift.  The occurring weights of $\mathcal U_n$
are the linear weights $e_i$, the cubic weights $2e_i+e_j$
($i\ne j$), and the weights $e_i+e_j+e_k$ with distinct indices;
the latter have multiplicity two.  Torus semistability implies
that $\mathbf1$ belongs to the rational span of the supported
weights.  Replace each supported linear weight $e_i$ by $3e_i$.
This preserves the rational span, and on a two-power diagonal
element it preserves the condition of exact invariance.
The resulting columns are exponent vectors of a cubic.

The whole support theorem ALPOGE\_GCT P-F30 applies to these
columns, including its arbitrary-lift-order and rank-deficient
cases.  It bounds the two-power determinant order by $M$ and,
at equality, forces exactly the cubic chain support
\[
                 2e_i+e_{i+1}\ (i<n),\qquad3e_n.
\]
The hook summand contains no cube weight.  Thus the last column
must have come from the linear coordinate $e_n$; all other
columns are the displayed hook weights, whose spaces are
one-dimensional.  Every such fixed semistable vector has all
$n$ coefficients nonzero and is diagonally equivalent to $z_n$.
This proves that every nonzero $\mu_M$-fixed quotient point
belongs to the same $GL_n/SL_n$ grading orbit.  Conversely that
orbit is fixed, by (K-024b).

The quotient is smooth at this orbit, because the closed orbit
has trivial $SL_n$ stabilizer.  Its tangent representation is
\[
                  \mathcal U_n/(\mathfrak{sl}_n\cdot z_n).
\]
The complete fixed space of $g$ in $\mathcal U_n$ has dimension
$n$ by the same chain-support classification.  More explicitly,
adjoin any hypothetical additional fixed weight to the known
chain support; the rational-span hypothesis still holds, and
P-F30 excludes that extra weight.  Meanwhile
$\dim(\mathfrak{sl}_n)^g=n-1$.  Thus its fixed tangent has
dimension one.  The closure of the nonzero grading orbit has
only the origin as boundary, since the coordinate grading is
positive.  Its stabilizer is exactly $\mu_M$ by the full
stabilizer calculation.  Its finitely generated degree semigroup
has group of differences $M\mathbb Z$, hence is cofinite after
division by $M$; its density is one.

Apply the unique-component Koszul formula P-F23.  At
$\xi=\det g=\zeta^{\sum_iw_i}$ the exact normalized residue is
\[
 C_n(\xi)=
 \frac{\displaystyle\prod_{i\ne j}(1-\zeta^{w_i-w_j})}
 {\displaystyle
  \prod_{w_i\ne0}(1-\zeta^{w_i})
  \prod_{\substack{i\ne j\\2w_i+w_j\ne0}}
                          (1-\zeta^{2w_i+w_j})
  \prod_{\substack{i<j<k\\w_i+w_j+w_k\ne0}}
                          (1-\zeta^{w_i+w_j+w_k})^2}.
                                                               \tag{K-024c}
\]
All congruences in this formula are modulo $M$.  The multiplicity
two in the last product is the actual distinct-letter weight
multiplicity of $\mathbb S_{(2,1)}W$.  Pullback is by the point
action, as in K-014 and K-019.
Every displayed factor is nonzero.  Moreover if $\eta$ is any
nontrivial two-power root of order $2^b$, then
\[
                 (1-\eta)\sum_{j=0}^{2^{b-1}-1}\eta^j=2.
\]
Thus every factor is a unit in $\mathbb Z[\zeta,1/2]$.
Since $\sum_iw_i$ is odd, this equals
$\mathbb Z[\xi,1/2]$.  The residue is a unit, proving the simple
pole and (K-024a).  Galois conjugation gives all primitive roots.

Finally let the two-part of $d$ be at least $M$ and $d\ne M$.
Then $\mu_d$ contains $\mu_M$.  Any nonzero $\mu_d$-fixed point
would lie on the just classified grading orbit, whose stabilizer
is exactly $\mu_M$, a contradiction.  Its fixed scheme is
therefore supported only at the origin.  The finite-length
Koszul argument of P-F23, also used for the pole cutoff in K-009,
shows that the global Hilbert series is regular at a primitive
$d$th root.  This proves the last assertion.
\end{proof}

The characteristic-zero smoothness argument is logically separate
from reduction modulo three of the integral Hilbert series.
In particular R-NEW-004 uses a coefficient system of determinant
three, and no characteristic-three smoothness or equality of
invariant rings after reduction is asserted or needed.
